\section{Wstęp}
Miara Lebesgue'a, czyli miara podzbiorów przestrzeni rzeczywistej $\R$ (i ogólniej, $\Rn$), pochodzi z pracy Henriego Lebesgue'a, która dotyczyła rozszerzenia pojęcia całki na znacznie szerszą klasę funkcji niż całka Riemanna (tzw. całka Lebesgue'a). Teoria miary niemal od początku stawiała pytanie, czy istnieje miara, czyli funkcja przypisująca podzbiorom ustalonego zbioru nieujemne wartości liczbowe lub $+\infty$ w taki sposób, że zbiór pusty ma miarę zero, a miara sumy przeliczalnej rodziny parami rozłącznych zbiorów jest równa sumie ich miar, która ponadto (1) jest określona na wszystkich podzbiorach $\R$, (2) jest niezmiennicza na przesunięcia oraz (3) przypisuje przedziałowi jednostkowemu wartość 1. W teorii zbiorów z aksjomatem wyboru taka miara nie istnieje. Miara Lebesgue'a spełnia drugi i trzeci z tych warunków, a zamiast pierwszego podaje kryterium, które dzieli wszystkie podzbiory zbioru liczb rzeczywistych na mierzalne oraz niemierzalne. Definicja miary Lebesgue'a $\lambda$ opiera się na definicji miary zewnętrznej Lebesgue'a $\lambda^{*}$, która jest określona na zbiorze wszystkich podzbiorów przestrzeni $\Rn$. Miara zewnętrzna Lebesgue'a jest niezmiennicza na przesunięcia oraz kostce jednostkowej przypisuje wartość 1, co znaczy, że jest odpowiednia do pomiaru wielkości geometrycznych, takich jak długość, pole powierzchni czy objętość. Nie jest ona jednak przeliczalnie addytywna, co oznacza, że miara sumy zbiorów rozłącznych może być mniejsza niż suma miar tych zbiorów. Właściwa miara Lebesgue'a nie ma już tej wady (tzn.\ jest przeliczalnie addytywna): jest równa mierze zewnętrznej Lebesgue'a, ale jej dziedzina jest ograniczona do zbiorów spełniających warunek Carathéodory'ego; zbiory te tworzą $\sigma$-ciało, które jest dziedziną miary Lebesgue'a. W pracy tej odtwarza się krok po kroku proces wiodący do konstrukcji tej miary.
